% Section 3: Vulnérabilités identifiées
% Fichier: sections/03_vulnerabilities.tex

\chapter{Vulnérabilités identifiées}

\section{Challenges ⭐ (1 étoile)}

\begin{vulnerability}{Score Board ⭐}{A05:2021 - Security Misconfiguration}{⭐}
    \subsection*{Description}
    Trouver la page Score Board soigneusement cachée.
    
    \subsection*{Technique d'exploitation}
    1. Clic droit sur la page d'accueil > Inspecter > Sources
    2. Rechercher `path` ou `score` dans `main.js`
    3. Accéder à l'URL : `http://[ip]:3000/#/score-board`
    
    \subsection*{Impact sur la sécurité}
    \begin{itemize}
        \item \textbf{Sévérité:} Basse
        \item \textbf{Impact:} Exposition de la liste des vulnérabilités
    \end{itemize}
    
    \subsection*{Recommandations de correction}
    \begin{recommendation}{Moyenne}
        Restreindre l'accès au score board aux seuls administrateurs.
    \end{recommendation}
\end{vulnerability}

\begin{vulnerability}{DOM XSS ⭐}{A03:2021 - Injection}{⭐}
    \subsection*{Description}
    Réaliser une attaque DOM XSS avec un payload iframe.
    
    \subsection*{Technique d'exploitation}
    \begin{payload}
<iframe src="javascript:alert('xss')">
    \end{payload}
    1. Cliquer sur l'icône de recherche
    2. Saisir le payload dans la barre de recherche
    3. Appuyer sur Entrée
    
    \subsection*{Impact sur la sécurité}
    \begin{itemize}
        \item \textbf{Sévérité:} Moyenne
        \item \textbf{Impact:} Exécution de code JavaScript dans le navigateur de la victime
    \end{itemize}
    
    \subsection*{Recommandations de correction}
    \begin{recommendation}{Critique}
        Échapper toutes les données utilisateur avant de les insérer dans le DOM.
    \end{recommendation}
\end{vulnerability}

\begin{vulnerability}{Confidential Document ⭐}{A01:2021 - Broken Access Control}{⭐}
    \subsection*{Description}
    Accéder à un document confidentiel via path traversal.
    
    \subsection*{Technique d'exploitation}
    1. Aller sur la page "About Us"
    2. Cliquer sur le lien des conditions d'utilisation
    3. Supprimer "legal.md" dans l'URL pour accéder au dossier ftp
    4. Ouvrir "acquisitions.md"
    
    \subsection*{Impact sur la sécurité}
    \begin{itemize}
        \item \textbf{Sévérité:} Moyenne
        \item \textbf{Impact:} Exposition de données sensibles
    \end{itemize}
    
    \subsection*{Recommandations de correction}
    \begin{recommendation}{Élevée}
        Valider et restreindre strictement l'accès aux fichiers.
    \end{recommendation}
\end{vulnerability}

\section{Challenges ⭐⭐ (2 étoiles)}

\begin{vulnerability}{Login Admin ⭐⭐}{A03:2021 - Injection}{⭐⭐}
    \subsection*{Description}
    Se connecter avec le compte administrateur via injection SQL.
    
    \subsection*{Technique d'exploitation}
    \begin{payload}
Email: ' OR TRUE --
Password: [n'importe quelle valeur]
    \end{payload}
    
    \subsection*{Impact sur la sécurité}
    \begin{itemize}
        \item \textbf{Sévérité:} Élevée
        \item \textbf{Impact:} Accès administrateur complet
    \end{itemize}
    
    \subsection*{Recommandations de correction}
    \begin{recommendation}{Critique}
        Utiliser des requêtes paramétrées.
    \end{recommendation}
\end{vulnerability}

\begin{vulnerability}{Reflected XSS ⭐⭐}{A03:2021 - Injection}{⭐⭐}
    \subsection*{Description}
    Réaliser une attaque XSS réfléchie dans le suivi de commande.
    
    \subsection*{Technique d'exploitation}
    \begin{payload}
<iframe src="javascript:alert(`xss`)">
    \end{payload}
    1. Passer une commande
    2. Aller dans l'historique des commandes
    3. Cliquer sur "Track Order"
    4. Remplacer l'ID de commande par le payload URL encodé
    
    \subsection*{Impact sur la sécurité}
    \begin{itemize}
        \item \textbf{Sévérité:} Moyenne
        \item \textbf{Impact:} Exécution de code dans le navigateur
    \end{itemize}
    
    \subsection*{Recommandations de correction}
    \begin{recommendation}{Critique}
        Échapper les données et implémenter une CSP.
    \end{recommendation}
\end{vulnerability}

\section{Challenges ⭐⭐⭐ (3 étoiles)}

\begin{vulnerability}{Login Bender ⭐⭐⭐}{A03:2021 - Injection}{⭐⭐⭐}
    \subsection*{Description}
    Se connecter avec le compte de Bender.
    
    \subsection*{Technique d'exploitation}
    \begin{payload}
Email: bender@juice-sh.op' --
Password: [n'importe quelle valeur]
    \end{payload}
    
    \subsection*{Impact sur la sécurité}
    \begin{itemize}
        \item \textbf{Sévérité:} Élevée
        \item \textbf{Impact:} Compromission du compte de Bender
    \end{itemize}
    
    \subsection*{Recommandations de correction}
    \begin{recommendation}{Critique}
        Utiliser des requêtes paramétrées.
    \end{recommendation}
\end{vulnerability}

\begin{vulnerability}{CSRF ⭐⭐⭐}{A01:2021 - Broken Access Control}{⭐⭐⭐}
    \subsection*{Description}
    Changer le nom d'un utilisateur via une attaque CSRF.
    
    \subsection*{Technique d'exploitation}
    \begin{payload}
<form action="http://[ip]:3000/profile" method="POST">
  <input name="username" value="CSRF"/>
  <input type="submit"/>
</form>
<script>document.forms[0].submit();</script>
    \end{payload}
    
    \subsection*{Impact sur la sécurité}
    \begin{itemize}
        \item \textbf{Sévérité:} Élevée
        \item \textbf{Impact:} Modification non autorisée de données
    \end{itemize}
    
    \subsection*{Recommandations de correction}
    \begin{recommendation}{Critique}
        Utiliser des tokens CSRF et configurer SameSite=Lax pour les cookies.
    \end{recommendation}
\end{vulnerability}

\section{Challenges ⭐⭐⭐⭐ (4 étoiles)}

\begin{vulnerability}{NoSQL Injection ⭐⭐⭐⭐}{A03:2021 - Injection}{⭐⭐⭐⭐}
    \subsection*{Description}
    Exploiter une injection NoSQL pour manipuler la base de données.
    
    \subsection*{Technique d'exploitation}
    Utiliser des opérateurs NoSQL comme `$ne` (not equal) pour contourner les contrôles.
    
    \subsection*{Impact sur la sécurité}
    \begin{itemize}
        \item \textbf{Sévérité:} Élevée
        \item \textbf{Impact:} Manipulation de données, contournement d'authentification
    \end{itemize}
    
    \subsection*{Recommandations de correction}
    \begin{recommendation}{Critique}
        Valider et nettoyer toutes les entrées utilisateur pour les requêtes NoSQL.
    \end{recommendation}
\end{vulnerability}

\section{Challenges ⭐⭐⭐⭐⭐ (5 étoiles)}

\begin{vulnerability}{Unsigned JWT ⭐⭐⭐⭐⭐}{A02:2021 - Cryptographic Failures}{⭐⭐⭐⭐⭐}
    \subsection*{Description}
    Forger un token JWT non signé pour usurper l'identité d'un utilisateur.
    
    \subsection*{Technique d'exploitation}
    1. Se connecter et récupérer le token JWT
    2. Décoder le token
    3. Changer l'email en `jwtn3d@juice-sh.op`
    4. Changer l'algorithme (`alg`) de `RS256` à `none`
    5. Réencoder le token et l'utiliser
    
    \subsection*{Impact sur la sécurité}
    \begin{itemize}
        \item \textbf{Sévérité:} Critique
        \item \textbf{Impact:} Usurpation d'identité complète
    \end{itemize}
    
    \subsection*{Recommandations de correction}
    \begin{recommendation}{Critique}
        Vérifier systématiquement la signature des JWT et ne jamais accepter l'algorithme `none`.
    \end{recommendation}
\end{vulnerability}

\section{Challenges ⭐⭐⭐⭐⭐⭐ (6 étoiles)}

\begin{vulnerability}{SSRF ⭐⭐⭐⭐⭐⭐}{A10:2021 - Server-Side Request Forgery}{⭐⭐⭐⭐⭐⭐}
    \subsection*{Description}
    Demander une ressource cachée sur le serveur via SSRF.
    
    \subsection*{Technique d'exploitation}
    1. Utiliser le champ "Image URL" sur la page de profil
    2. Saisir l'URL de la ressource locale : `http://localhost:3000/solve/challenges/server-side?key=tRy_H4rd3r_n0thIng_iS_Imp0ssibl3`
    3. Cliquer sur "Link Image"
    
    \subsection*{Impact sur la sécurité}
    \begin{itemize}
        \item \textbf{Sévérité:} Critique
        \item \textbf{Impact:} Accès à des ressources internes du serveur
    \end{itemize}
    
    \subsection*{Recommandations de correction}
    \begin{recommendation}{Critique}
        Restreindre les requêtes sortantes du serveur et valider les URLs utilisateur.
    \end{recommendation}
\end{vulnerability}

\begin{vulnerability}{Successful RCE DoS ⭐⭐⭐⭐⭐⭐}{A08:2021 - Software and Data Integrity Failures}{⭐⭐⭐⭐⭐⭐}
    \subsection*{Description}
    Réaliser une exécution de code à distance (RCE) via désérialisation non sécurisée.
    
    \subsection*{Impact sur la sécurité}
    \begin{itemize}
        \item \textbf{Sévérité:} Critique
        \item \textbf{Impact:} Contrôle complet du serveur
    \end{itemize}
    
    \subsection*{Recommandations de correction}
    \begin{recommendation}{Critique}
        Éviter la désérialisation de données non fiables et utiliser des formats de données sûrs.
    \end{recommendation}
\end{vulnerability}
